\medskip

\emph{Cet exercice est un questionnaire à choix multiples. Pour chacune des questions suivantes, une seule des quatre réponses proposées est exacte. Une réponse fausse, une réponse multiple ou l'absence de réponse à une question ne rapporte ni n'enlève de point.\\
Pour répondre, indiquer sur la copie le numéro de la question et la lettre de la réponse choisie.\\Aucune justification n'est demandée.}

\medskip

Pour les questions 1 à 3 ci-dessous, on considère une fonction $f$ définie et deux fois dérivable sur $\R$. La courbe de sa fonction dérivée $f'$ est donnée ci-dessous.

On admet que $f'$ admet un maximum en $- \dfrac{3}{2}$ et que sa courbe coupe l'axe des abscisses au point de coordonnées $\left(- \dfrac12~;~0\right)$.

\begin{minipage}{0.5\linewidth}
\textbf{Question 1 }:

\textbf{a.~~} La fonction $f$ admet un maximum en $- \dfrac{3}{2}$ ;

\textbf{b.~~}  La fonction $f$ admet un maximum en $- \dfrac{1}{2}$ ;

\textbf{c.~~} La fonction $f$ admet un minimum en $- \dfrac{1}{2}$; 

\textbf{d.~~}  Au point d'abscisse $-1$, la courbe de la
fonction $f$ admet une tangente horizontale.
\end{minipage}\hfill
\begin{minipage}{0.48\linewidth}
\begin{center}
On rappelle que la courbe ci-dessous représente la fonction dérivée $f'$ de $f$.
\end{center}

\psset{unit=1.2cm}
\begin{pspicture*}(-5.2,-2.6)(1,1)
\psgrid[gridlabels=0pt,subgriddiv=4,gridwidth=0.25pt,subgridwidth=0.15pt]
\psaxes[linewidth=1.25pt,labelFontSize=\scriptstyle]{->}(0,0)(-5.2,-2.6)(1,1)

\psplot[plotpoints=2000,linewidth=1.25pt,linecolor=red]{-5}{3}{2 x mul 1 add 2.71828 x exp mul neg}
\end{pspicture*}
\end{minipage}
Par lecture graphique, la dérivée est positive sur $\left]- \infty~;~- \dfrac12\right[$ et négative sur $\left[- \dfrac12~;~+ \infty\right[$ ; $f$ est donc croissante puis décroissante : elle admet donc un maximum en $- \dfrac12$. Réponse \textbf{b.}
\medskip

\textbf{Question 2 }:

\begin{center}
\begin{tabularx}{\linewidth}{X X}
\textbf{a.~~}La fonction $f$ est convexe sur $\left]- \infty~;~- \dfrac32\right[$ ;&
\textbf{b.~~}La fonction $f$ est convexe sur $\left]- \infty;~- \dfrac12\right[$ ;\\
\textbf{c.~~}La courbe $\mathcal{C}_f$ représentant la fonction $f$ n'admet pas de point d'inflexion ; &\textbf{d.~~} la fonction $f$ est concave sur $\left] - \infty~;~- \dfrac12\right[$.
\end{tabularx}
\end{center}

$f'$ est croissante sur $\left]- \infty~;~- \dfrac32\right[$ : $f$ est convexe sur cet intervalle. Réponse \textbf{a.}
\medskip

\textbf{Question 3}:

La dérivée seconde $f''$  de la fonction $f$ vérifie :

\begin{center}
\begin{tabularx}{\linewidth}{X X}
\textbf{a.~~} $f''(x) \geqslant  0$ pour $x \in \left]-\infty~;~- \dfrac12\right[$ ; &\textbf{b.~~} $f''(x) \geqslant  0$ pour $x \in [- 2~;~- 1]$ ;\\
\textbf{c.~~} $f''\left(- \dfrac32 \right) = 0$ ;&\textbf{d.~~} $f''(- 3) = 0$.
\end{tabularx}
\end{center}
On a déjà vu que $f'$ est croissante sur $\left]- \infty~;~- \dfrac32\right[$ ; elle est décroissante sur $\left]- \dfrac32~;~+ \infty\right[$, donc 

$f''\left(- \dfrac32\right) = 0$. Réponse \textbf{c.}

\medskip

\textbf{Question 4 :}

\medskip

On considère trois suites $\left(u_n\right)$, $\left(v_n\right)$ et $\left(w_n\right)$. On sait que, pour tout entier naturel $n$, on a : $u_n \leqslant v_n\leqslant  w_n$ et de plus: $\displaystyle\lim_{n \to + \infty} u_n= 1$ et $\displaystyle\lim_{n \to + \infty} w_n= 3$.

On peut alors affirmer que :

\begin{center}
\begin{tabularx}{\linewidth}{X X}
\textbf{a.~~} la suite $\left(v_n\right)$ converge ;&\textbf{b.~~} Si la suite 
$\left(u_n\right)$ est croissante alors la suite $\left(v_n\right)$ est minorée par $u_0$ ;\\
\textbf{c.~~} $1 \leqslant  v_0 \leqslant 3$ ;&\textbf{d.~~} la suite $\left(v_n\right)$ diverge.
\end{tabularx}
\end{center}

Si $\left(u_n\right)$ est croissante, on a pour tout $n \in \N$,\: $u_0 \leqslant u_n \leqslant v_n$, donc la suite $\left(v_n\right)$ est minorée par $u_0$. Réponse \textbf{b.}

\medskip

\textbf{Question 5}:

\medskip

On considère une suite $\left(u_n\right)$ telle que, pour tout entier naturel $n$ non nul: $u_n \leqslant  u_{n+1}  \leqslant \dfrac1n$.

On peut alors affirmer que : 

\begin{center}
\begin{tabularx}{\linewidth}{X X}
\textbf{a.~~}la suite $\left(u_n\right)$ diverge ;&\textbf{b.~~}la suite $\left(u_n\right)$ converge ;\\
\textbf{c.~~}$\displaystyle\lim_{n \to + \infty} u_n =  0$ ;&\textbf{d.~~} $\displaystyle\lim_{n \to + \infty} u_n =  1$.
\end{tabularx}
\end{center}

Pour tout entier $n \in \N^{*}$, \: $u_n \leqslant u_{n+1} \leqslant \dfrac1n$ : la suite 
$\left(u_n\right)$ est donc croissante.

D'autre part : pour $n \ne 0$, on peut écrire : $u_n \leqslant u_{n+1} \leqslant \dfrac1n \leqslant 1$.

Conclusion  : la suite $\left(u_n\right)$ est croissante et elle majorée par 1 : elle converge donc. Réponse \textbf{b.}
\medskip

\textbf{Question 6}:

\medskip

On considère $\left(u_n\right)$ une suite réelle telle que pour tout entier naturel $n$, on a : $n < u_n < n + 1$. 

On peut affirmer que:

\begin{center}
\begin{tabularx}{\linewidth}{X X}
\textbf{a.~~}Il existe un entier naturel $N$ tel que $u_N$ est un entier ;&\textbf{b.~~} la suite $\left(u_n\right)$ est croissante ;\\
\textbf{c.~~} la suite $\left(u_n\right)$ est convergente ;&\textbf{d.~~}La suite $\left(u_n\right)$ n'a pas de limite.
\end{tabularx}
\end{center}

On a au rang $n$ : $n < u_n < n+1$ et au rang $n+1$ :

$n + 1 < u_{n+1} < u_{n+2}$, donc $u_n < n+1 < u_{n+1}$ : la suite $\left(u_n\right)$ est croissante. Réponse \textbf{b.}
